\\documentclass[12pt,a4paper]{article}
\\usepackage[a4paper,margin=1in]{geometry}
\\usepackage{times}
\\usepackage{setspace}
\\usepackage{booktabs}
\\usepackage{amsmath,amssymb}
\\usepackage{xcolor}
\\usepackage{listings}
\\usepackage{hyperref}
\\usepackage{enumitem}
\\usepackage{fancyhdr}
\\usepackage{titlesec}

\\hypersetup{colorlinks=true,linkcolor=blue,citecolor=blue,urlcolor=blue,pdftitle={Email Spam Detection Using Machine Learning}}
\\onehalfspacing
\\pagestyle{fancy}
\\fancyhf{}
\\fancyhead[L]{Email Spam Detection}
\\fancyhead[R]{Reg. No.: 24bcb7272}
\\fancyfoot[C]{\\thepage}

\\definecolor{codegreen}{rgb}{0,0.5,0}
\\definecolor{codegray}{rgb}{0.5,0.5,0.5}
\\definecolor{codepurple}{rgb}{0.58,0,0.82}
\\lstset{language=Python,basicstyle=\\ttfamily\\small,keywordstyle=\\color{blue},commentstyle=\\color{codegreen},stringstyle=\\color{codepurple},numbers=left,numberstyle=\\tiny\\color{codegray},stepnumber=1,numbersep=8pt,frame=single,breaklines=true,showstringspaces=false}
\\titleformat{\\section}{\\large\\bfseries}{\\thesection.}{0.5em}{}
\\titleformat{\\subsection}{\\normalsize\\bfseries}{\\thesubsection}{0.5em}{}

\\title{Email Spam Detection Using Machine Learning}
\\author{Reg. No.: 24bcb7272}
\\date{Academic Year 2026--2027}

\\begin{document}
\\begin{titlepage}
\\centering
\\vspace*{1cm}
{\\Large\\bfseries VIT-AP UNIVERSITY\\par}
\\vspace{0.3cm}
{\\large Amaravati, Andhra Pradesh, India\\par}
\\vspace{1.2cm}
{\\Large\\bfseries PROJECT REPORT\\par}
\\vspace{0.5cm}
{\\large Subject: Machine Learning / Data Mining\\par}
{\\large Course Code: CSEXXXX\\par}
\\vspace{1.2cm}
{\\LARGE\\bfseries Email Spam Detection Using Machine Learning\\par}
\\vspace{1.5cm}
{\\large Submitted by\\par}
\\vspace{0.4cm}
{\\bfseries Regd. No.: 24bcb7272\\par}
\\vspace{1.2cm}
{\\large Under the Guidance of\\par}
\\vspace{0.3cm}
{\\bfseries Dr. [Faculty Name]\\par}
\\vfill
{\\large School of Computer Science and Engineering (SCOPE)\\\\ VIT-AP University\\\\ Amaravati, Andhra Pradesh, India\\par}
\\vspace{0.5cm}
{\\large Academic Year 2026--2027\\par}
\\end{titlepage}

\\section*{Abstract}
Spam detection is a text-classification problem in which messages are classified as unsolicited (spam) or legitimate (ham). This project implements a supervised machine-learning workflow using the SMS Spam Collection dataset. Messages are normalized, filtered using English stopwords, stemmed, and represented with a bag-of-words vectorizer. Random Forest, Decision Tree, and Multinomial Naive Bayes classifiers are trained and evaluated with a held-out test set. Accuracy, confusion matrices, and per-class precision, recall, and F1-score are reported. The evaluation is intended to show the trade-off between detecting spam and avoiding false alarms; measured scores should be generated by running the supplied code against the chosen dataset.

\\textbf{Keywords:} Spam detection, natural language processing, text classification, bag of words, Naive Bayes.

\\section{Introduction}
Email and messaging services are widely used in personal, academic, and professional communication. Unsolicited messages can waste time and may contain fraudulent offers or unsafe links. Automated filtering can help prioritize suspicious messages by learning patterns from examples labeled as spam or ham.

\\subsection{Problem Statement}
Manual filtering is inefficient at scale. The project explores whether message text can be transformed into numerical features and classified using supervised learning, while making false positives and missed spam visible through class-specific metrics.

\\subsection{Objectives}
\\begin{enumerate}[leftmargin=2em]
    \\item Load and validate a labeled message dataset.
    \\item Normalize text, remove English stopwords, and apply Porter stemming.
    \\item Create bag-of-words features using a vocabulary learned from training examples only.
    \\item Compare Random Forest, Decision Tree, and Multinomial Naive Bayes classifiers.
    \\item Evaluate the models and save the selected model with its vectorizer.
\\end{enumerate}

\\section{Dataset and Methodology}
The code expects a CSV named \\texttt{spam.csv} with \\texttt{v1} and \\texttt{v2} columns, representing labels and message text. Extra columns in the common SMS Spam Collection CSV are ignored. Labels are validated as \\texttt{ham} or \\texttt{spam}. The class distribution is imbalanced, so accuracy is interpreted alongside spam precision and recall.

\\subsection{Preprocessing}
Each message is converted to lowercase; nonalphabetic characters are replaced with spaces; English stopwords are removed; and the remaining tokens are reduced using the NLTK Porter stemmer. Empty messages are handled as empty strings.

\\subsection{Feature Extraction and Data Split}
A \\texttt{CountVectorizer} with at most 4,000 features converts text into token-count vectors. The dataset is split into training and test partitions using a fixed random seed and stratification by label. The vectorizer is fitted on training text only, then applied to test text, preventing information from the test set leaking into the learned vocabulary.

\\subsection{Classifiers and Evaluation}
The project compares Random Forest, Decision Tree, and Multinomial Naive Bayes. Confusion matrices use ham as class 0 and spam as class 1. Accuracy, precision, recall, and F1-score are printed for each model. In particular, spam recall measures the fraction of spam examples found, while ham precision / spam false positives help indicate the risk of misclassifying legitimate messages.

\\section{Results and Discussion}
The Python program and notebook calculate evaluation metrics when run on the dataset. Results are not hard-coded here because they can differ depending on the dataset file and software versions. Insert the actual output from your run in this section before submission.

\\begin{table}[h]
\\centering
\\begin{tabular}{lrrrr}
\\toprule
Model & Accuracy & Spam precision & Spam recall & Spam F1 \\\\
\\midrule
Random Forest & -- & -- & -- & -- \\\\
Decision Tree & -- & -- & -- & -- \\\\
Multinomial Naive Bayes & -- & -- & -- & -- \\\\
\\bottomrule
\\end{tabular}
\\caption{Complete this table with the classification report produced by the script.}
\\label{tab:metrics}
\\end{table}

The class distribution is not balanced, so the majority ham class can make accuracy appear high even when some spam is missed. Model selection should consider the application’s tolerance for false positives and missed spam, rather than relying on a single metric.

\\subsection{Limitations}
\\begin{itemize}
    \\item Results depend on the dataset’s coverage, label quality, and message style.
    \\item Bag-of-words features do not represent word order or message metadata.
    \\item Real-world spam evolves; the model should be tested on recent, representative data before deployment.
    \\item A false positive may hide a legitimate message, so a model prediction should not be treated as definitive proof.
\\end{itemize}

\\section{Conclusion}
This project demonstrates a reproducible baseline for spam classification using text preprocessing, bag-of-words features, and supervised classifiers. The supplied notebook and Python script report multiple evaluation metrics and save the Multinomial Naive Bayes model together with the vectorizer. Future work could assess TF--IDF features, calibrated decision thresholds, newer datasets, and message metadata.

\\section*{References}
\\begin{enumerate}[leftmargin=2em]
    \\item SMS Spam Collection dataset, UCI Machine Learning Repository / Kaggle mirror: \\url{https://www.kaggle.com/datasets/uciml/sms-spam-collection-dataset}.
    \\item Bird, S., Klein, E., and Loper, E. \\textit{Natural Language Processing with Python}. O'Reilly Media.
    \\item Pedregosa, F. et al. Scikit-learn: Machine Learning in Python. \\textit{Journal of Machine Learning Research}, 12, 2011.
\\end{enumerate}
\\end{document}
